\documentclass[11pt,a4paper]{article}
\usepackage[italian]{babel}
\usepackage[margin=2.5cm]{geometry}
\usepackage{parskip}
\usepackage{amsmath, amssymb, amsthm}
\usepackage{booktabs}
\usepackage{array}
\usepackage{xcolor}
\usepackage{needspace}
\usepackage{enumitem}
\usepackage{listings}
\usepackage{tikz}
\usepackage[most]{tcolorbox}
\usepackage[hidelinks]{hyperref}
\usetikzlibrary{decorations.pathreplacing, shapes.geometric, arrows.meta, positioning}

% Niente vedove/orfane: i paragrafi non lasciano righe isolate a fine/inizio pagina
\widowpenalty=10000
\clubpenalty=10000

% Elenchi più compatti
\setlist{itemsep=2pt, parsep=0pt, topsep=4pt, partopsep=0pt}

% --- Riquadri con bordo nero, senza numero (si spezzano tra due pagine) ---
% Uso: \begin{definizione} ... \end{definizione}
%      \begin{teorema}[Nome] ... \end{teorema}  (il nome è facoltativo)
\tcbset{nota/.style={enhanced jigsaw, breakable, colback=white,
  colframe=black, boxrule=0.5pt, sharp corners}}
\NewTColorBox{definizione}{}{nota,
  before upper={\textbf{Definizione.}\ }}
\NewTColorBox{teorema}{o}{nota,
  before upper={\textbf{Teorema\IfValueT{#1}{ (#1)}.}\ }}
\NewTColorBox{proposizione}{o}{nota,
  before upper={\textbf{Proposizione\IfValueT{#1}{ (#1)}.}\ }}
\NewTColorBox{proprieta}{o}{nota,
  before upper={\textbf{Proprietà\IfValueT{#1}{ (#1)}.}\ }}
\NewTColorBox{assioma}{o}{nota, boxrule=1pt,
  before upper={\textbf{Assioma\IfValueT{#1}{ (#1)}.}\ }}

% --- Ambienti semplici (senza riquadro) ---
% Il titolo sta in cima al blocco, anche se il contenuto inizia con un riquadro o
% con codice; e non resta mai da solo in fondo a una pagina.
% Uso: \begin{esempio}[Titolo facoltativo] ... \end{esempio}
\newcommand{\newrunin}[2]{%
  \NewDocumentEnvironment{#1}{o}{%
    \par\Needspace{4\baselineskip}\addvspace{\medskipamount}\noindent
    \textbf{#2}\IfValueT{##1}{ (##1)}\textbf{.}\ \ignorespaces}%
  {\par\addvspace{\medskipamount}}}
\newrunin{esempio}{Esempio}
\newrunin{esercizio}{Esercizio}
\NewTColorBox{osservazione}{}{nota,
  before upper={\textbf{Osservazione.}\ }}

% V e F nelle tavole di verità
\newcommand{\V}{\textsf{V}}
\newcommand{\F}{\textsf{F}}

% --- Codice C ---
% Uso: \begin{codice} ... \end{codice}      codice C numerato
%      \begin{comando} ... \end{comando}    comandi da terminale
%      \begin{risultato} ... \end{risultato} output di un programma
%      \lstinline|printf("ciao");|           codice dentro al testo
\lstdefinestyle{c}{
  language=C,
  basicstyle=\ttfamily\small,
  keywordstyle=\color{blue}\bfseries,
  commentstyle=\color{gray}\itshape,
  stringstyle=\color{red!70!black},
  numbers=left,
  numberstyle=\tiny\color{gray},
  numbersep=8pt,
  columns=flexible,
  keepspaces=true,
  breaklines=true,
  showstringspaces=false,
  tabsize=4,
  morekeywords={include, define},
  % lettere accentate nei commenti
  literate={à}{{\`a}}1 {è}{{\`e}}1 {é}{{\'e}}1 {ì}{{\`i}}1
           {ò}{{\`o}}1 {ù}{{\`u}}1 {È}{{\`E}}1 {À}{{\`A}}1
}
\lstset{style=c}
\newtcblisting{codice}{listing only, nota, left=6mm,
  listing options={style=c}}
\newtcblisting{comando}{listing only, nota,
  listing options={basicstyle=\ttfamily\small, numbers=none, language={},
    columns=flexible, keepspaces=true, breaklines=true}}
\newtcblisting{risultato}{listing only, enhanced jigsaw, breakable,
  colback=black!5, colframe=black, boxrule=0.5pt, sharp corners,
  listing options={basicstyle=\ttfamily\small, numbers=none, language={},
    columns=flexible, keepspaces=true, breaklines=true}}

% --- Diagrammi di flusso (simboli standard) ---
\tikzset{
  terminale/.style = {ellipse, draw, minimum width=2.5cm, minimum height=0.9cm},
  io/.style        = {trapezium, trapezium left angle=70, trapezium right angle=110,
                      draw, minimum width=2.5cm, minimum height=0.9cm},
  processo/.style  = {rectangle, draw, minimum width=2.5cm, minimum height=0.9cm},
  decisione/.style = {diamond, draw, aspect=2, minimum width=2.5cm},
  freccia/.style   = {thick, -{Stealth}}
}

\title{Progettazione di reti logiche combinatorie}
\author{Marco Cerbella}
\date{Lezione 4}

\begin{document}
\maketitle

Una rete logica \emph{combinatoria} ha un comportamento che dipende solo dagli ingressi, perché è priva di memoria. Vediamo come si progetta e si ottimizza.

\section{Le fasi della progettazione}

Progettare un circuito combinatorio comporta quattro passi successivi.

\begin{enumerate}
    \item \textbf{Comprensione e analisi della specifica semi-formale.} La descrizione iniziale di un problema spesso non è completa e va riscritta per eliminare ogni ambiguità. Casi tipici: gestione delle condizioni limite o dei casi di errore, analisi dell'ambiente in cui la rete opera, scelta della codifica delle informazioni. Alla fine si ha ancora una specifica semi-formale, ma completa.
    \item \textbf{Formalizzazione della specifica.} Si traduce la specifica in un linguaggio formale, più adatto alla manipolazione: le \emph{tabelle di verità}, che specificano la relazione richiesta tra ingressi e uscite.
    \item \textbf{Sintesi.} Si passa dalle tabelle di verità a espressioni algebriche booleane. In questa fase si trascura ogni considerazione sulla qualità.
    \item \textbf{Ottimizzazione.} Fissato un criterio di qualità, si manipola l'espressione ottenuta nella sintesi per avere un'espressione funzionalmente equivalente ma migliore.
\end{enumerate}

Le caratteristiche che misurano la qualità di una rete sono generalmente l'area della realizzazione fisica, il tempo necessario a produrre un risultato, la potenza e l'energia assorbita (\emph{metriche d'area}, \emph{metriche di tempo}, \dots). Alcune metriche d'area comuni sono il numero di porte logiche generiche, il numero di porte a due ingressi, il numero di implicanti/implicati e il numero di letterali.

\section{Esempio: riconoscere pari e dispari}

\emph{Si progetti una rete combinatoria in grado di discriminare i numeri pari dai numeri dispari.}

\subsection*{Fase 1: specifica}

\begin{itemize}
    \item Ci interessano i numeri interi e dobbiamo fissare l'intervallo, in base all'applicazione. Scegliamo per semplicità i giorni della settimana, cioè l'intervallo $[0, 6]$.
    \item Servono $\lceil \log_2 7 \rceil = 3$ bit per rappresentarlo.
    \item L'uscita è binaria (pari o dispari), quindi basta un bit. Fissiamo la codifica $1 = $ pari e $0 = $ dispari.
    \item Con l'intervallo $[0,6]$ l'ingresso $111$ (il numero $7$) non è valido. Introduciamo un secondo bit di uscita per segnalare l'errore.
\end{itemize}

La specifica semi-formale completa è: \emph{una rete combinatoria con un ingresso a tre bit $X = \{x_0, x_1, x_2\}$ che rappresenta un intero in $[0,6]$ in codifica binaria naturale e due uscite di un bit ciascuna: $z$ indica se il valore è pari ($z = 1$) o dispari ($z = 0$), $e$ indica un errore ($e = 1$) o il funzionamento corretto ($e = 0$).}

\subsection*{Fase 2: formalizzazione}

La rete realizza le funzioni $z(x_0,x_1,x_2)$ ed $e(x_0,x_1,x_2)$, con $x_0$ bit più significativo:
\begin{center}
\begin{tabular}{ccc|cc}
\toprule
$x_0$ & $x_1$ & $x_2$ & $z$ & $e$ \\
\midrule
0 & 0 & 0 & 1 & 0 \\
0 & 0 & 1 & 0 & 0 \\
0 & 1 & 0 & 1 & 0 \\
0 & 1 & 1 & 0 & 0 \\
1 & 0 & 0 & 1 & 0 \\
1 & 0 & 1 & 0 & 0 \\
1 & 1 & 0 & 1 & 0 \\
1 & 1 & 1 & 0 & 1 \\
\bottomrule
\end{tabular}
\end{center}

\subsection*{Fase 3: sintesi}

Dalle tabelle si ricavano le espressioni in prima forma canonica:
\begin{align*}
z &= x_0'x_1'x_2' + x_0'x_1x_2' + x_0x_1'x_2' + x_0x_1x_2', \\
e &= x_0x_1x_2.
\end{align*}
Da qui si potrebbe già derivare un circuito, ma le proprietà dell'algebra booleana permettono di ottimizzare (in numero di letterali).

\subsection*{Fase 4: ottimizzazione}

Raccogliendo $x_2'$ e usando $y + y' = 1$ due volte:
\[
z = x_2'\,(x_0'x_1' + x_0'x_1 + x_0x_1' + x_0x_1) = x_2'\,\bigl(x_0'(x_1'+x_1) + x_0(x_1'+x_1)\bigr) = x_2'.
\]
Quindi $z = x_2'$: un numero è pari se e solo se l'ultimo bit è $0$. Le tecniche sistematiche di ottimizzazione sono oggetto delle lezioni successive (mappe di Karnaugh).

\section{Secondo esempio: display a sette segmenti}

Un display a $7$ segmenti è costituito da $7$ segmenti luminosi che si accendono o spengono a formare il simbolo di una cifra; per controllarlo servono $7$ segnali binari. Si progetti una rete combinatoria che gestisca un display a sette segmenti di una cifra decimale nell'intervallo $[0, 7]$ in codifica binaria naturale: la rete ha $3$ ingressi e $7$ uscite, una per segmento.

\medskip
\noindent\textbf{Riferimenti.} Libro di testo (C. Bolchini et al., \emph{Reti Logiche}, Maggioli): cap.\ 4, sottosezioni 4.1, 4.2, 4.3. Esercizi per casa: 4.1, 4.2, 4.5, 4.6, 4.7, 4.9.

\end{document}
